% ===========================================================================
\section{Gröbner Bases}\label{sec:groebner}
\warmth{0}\quad\whisper{The engine. Gröbner is to ideals what Gaussian elimination is to matrices.}

\begin{strand}
An ideal has many generating sets and no canonical basis. \keyword{Gröbner bases} fix this:
they are the right generating set once you pick a \emph{monomial order}. Think of them as the
Euclidean algorithm for $n\ge2$ variables, or Gaussian elimination for degree $\ge2$. The book's
creed: \keyword{nonlinear algebra is the next step after linear algebra}.
\end{strand}

\block{Monomial orders and initial terms}

\begin{definition}[monomial order, Def.\ 1.10]\label{def:order}
A total order $\prec$ on $\N^n$ (monomials) is a \keyword{monomial order} if
$(0,\dots,0)\preceq\mathbf a$ for all $\mathbf a$, and $\mathbf a\preceq\mathbf b\Rightarrow
\mathbf a+\mathbf c\preceq\mathbf b+\mathbf c$. Three standards (all with $x_1\succ\cdots\succ x_n$):
\emph{lex}, \emph{deglex} (degree first, then lex), \emph{degrevlex} (degree first, then
\emph{smallest} last variable wins). Any monomial order refines divisibility:
$\xb^{\mathbf a}\mid\xb^{\mathbf b}\Rightarrow\mathbf a\preceq\mathbf b$.
\end{definition}

The \keyword{initial monomial} $\inn_\prec(f)$ is the $\prec$-largest monomial of $f$.

\begin{yourturn}
Let $f=x^2+xz^2+y^3$ with $x\succ y\succ z$. Fill the three initial monomials:
\[
  \inn_{\mathrm{lex}}(f)=\fillin[1.6cm],\qquad
  \inn_{\mathrm{deglex}}(f)=\fillin[1.6cm],\qquad
  \inn_{\mathrm{degrevlex}}(f)=\fillin[1.6cm].
\]
\recall{For degrevlex, why does $y^3$ beat $xz^2$? Both have degree 3 -- compare the \emph{last} variable.}
\end{yourturn}

\block{Finiteness: Dickson and the Basis Theorem}

\begin{theorem}[Dickson's Lemma, Thm.\ 1.8]\label{thm:dickson}
Any infinite subset of $\N^n$ contains a pair $\{\mathbf a,\mathbf b\}$ with $\mathbf a\le\mathbf b$
(coordinatewise). Equivalently: any nonempty $\mathcal M\subseteq\N^n$ has a \emph{finite} set of
coordinatewise-minimal elements.
\end{theorem}
\begin{proof}
Induction on $n$. For $n=1$, any two naturals are comparable. \TODO{induction step: slice $\mathcal M$ by the last coordinate $i$, set $\mathcal M_i\subseteq\N^{n-1}$; if some $\mathcal M_i$ is infinite use IH, else \dots}
\end{proof}

For an ideal $I$, the \keyword{initial ideal} is $\inn_\prec(I)=\gen{\inn_\prec(f):f\in I\setminus\{0\}}$,
a monomial ideal. A finite $\mathcal G\subseteq I$ with $\inn_\prec(I)=\gen{\inn_\prec(g):g\in\mathcal G}$
is a \keyword{Gröbner basis} (it exists, by Dickson).

\begin{theorem}[a Gröbner basis generates, Thm.\ 1.13]\label{thm:gb-gen}
If $\mathcal G$ is a Gröbner basis of $I$, then $I=\gen{\mathcal G}$.
\end{theorem}
\begin{proof}
Suppose not. Pick $f\in I\setminus\gen{\mathcal G}$ with $\inn_\prec(f)=\xb^{\mathbf b}$ minimal.
Some $g\in\mathcal G$ has $\inn_\prec(g)\mid\xb^{\mathbf b}$, say $\xb^{\mathbf b}=\xb^{\mathbf c}\inn_\prec(g)$.
\TODO{form $f-\xb^{\mathbf c}g$: it lies in $I$, has strictly smaller initial monomial, contradiction}
\end{proof}

\begin{corollary}[Hilbert Basis Theorem, Cor.\ 1.14]\label{cor:hbt}
Every ideal in $K[\xb]$ is \fillin[3.5cm].
\end{corollary}

\block{Reduced Gröbner bases and standard monomials}

\begin{definition}[reduced, Def.\ 1.15]
$\mathcal G$ is a \keyword{reduced} Gröbner basis if (a) each leading coefficient is $1$, and (b) for
distinct $g,h\in\mathcal G$, no monomial of $g$ is a multiple of $\inn_\prec(h)$.
\end{definition}

\begin{theorem}[uniqueness + a basis, Thm.\ 1.16--1.17]\label{thm:reduced}
Fix $I$ and $\prec$. Then $I$ has a \emph{unique} reduced Gröbner basis. Moreover the
\keyword{standard monomials} $\mathcal S_\prec(I)=\{$monomials $\notin\inn_\prec(I)\}$ form a
$K$-vector-space basis of the quotient $K[\xb]/I$.
\end{theorem}
\begin{proof}[Proof of the basis claim]
Independence: a nonzero $f$ with image $0$ in $K[\xb]/I$ lies in $I$, so $\inn_\prec(f)\in\inn_\prec(I)$
is \emph{not} standard. Spanning: \TODO{take a minimal $\xb^{\mathbf c}$ not in the span; it lies in $\inn_\prec(I)$, divide, contradiction}
\end{proof}

\begin{yourturn}
(Symmetric polynomials, from Example 1.18.) The elementary symmetric polynomials
$\mathcal F=\{x+y+z,\ xy+xz+yz,\ xyz\}$ in $\Q[x,y,z]$ have reduced Gröbner basis (lex)
$\mathcal G=\{\underline{x}+y+z,\ \underline{y^2}+yz+z^2,\ \underline{z^3}\}$.
\begin{itemize}[leftmargin=1.3em,itemsep=0pt]
  \item List the \fillin[1cm] standard monomials $\mathcal S_\prec(I)$ (they index a basis of $\Q[x,y,z]/I$).
  \item This quotient is the regular representation of which group? \fillin[2cm]
\end{itemize}
\recall{\# standard monomials $=\dim_K K[\xb]/I=$ \# solutions counted with multiplicity. Here $=|S_3|=6$.}
\end{yourturn}

\begin{strand}
Buchberger's algorithm computes the reduced Gröbner basis from any $\mathcal F$. Two warnings the
book stresses: the \keyword{choice of order matters enormously} (Example 1.19: two random quartics in
$\PP^3$ give $5$ generators of degree $\le7$ in degrevlex, but $150$ of degree $\le73$ in lex), and
\keyword{lex is the order for elimination} (Chapter 4). Run a computer algebra system as you read.
\end{strand}

\loose{Exercise 1 (book): an ideal in $K[\xb]$ is principal $\iff$ its reduced Gröbner basis is a singleton. (One direction is easy; the other uses uniqueness.)}
